% ---------------------------------------------------------------------
\subsection{The walking camera (\fn{WalkView::advance}, \fn{Raster})}
\label{sec:math-walk}
The page \ui{Begehen} (\S\ref{sec:walk}) is drawn by a software renderer of its own (or, the same triangles, by the
graphics card). Everything below is in the room's coordinates: $x$ to the right, $y$ to the front, $z$ up, metres; the
eye stands at the listener's place plus the walked offset, $\mathbf e = \mathbf p_L + (p_x, p_y, z_f)$, where $z_f$ is
the height of the feet above the floor.

\paragraph{The camera.} The view has a yaw $\psi$ (positive = turned to the left, like the azimuth) and a pitch
$\vartheta$ (up). Its three axes are
\begin{equation}
  \mathbf f = (-\sin\psi\cos\vartheta,\ \cos\psi\cos\vartheta,\ \sin\vartheta),\qquad
  \mathbf r = (\cos\psi,\ \sin\psi,\ 0),\qquad
  \mathbf u = \mathbf r\times\mathbf f ,
\end{equation}
so $\psi = 0$ looks along $+y$, the Saal's front, and $\psi = 90^\circ$ along $-x$, to the left, exactly as an
Ambisonics source at azimuth $90^\circ$ lies. A point $\mathbf p$ goes into view space by
$\mathbf v = ((\mathbf p-\mathbf e)\cdot\mathbf r,\ (\mathbf p-\mathbf e)\cdot\mathbf u,\ (\mathbf p-\mathbf e)\cdot\mathbf f)$
and onto the picture of $w\times h$ pixels by the pinhole projection
\begin{equation}
  s_x = \tfrac w2 + F\,\frac{v_x}{v_z},\qquad s_y = \tfrac h2 - F\,\frac{v_y}{v_z},\qquad
  F = \frac{w/2}{\tan(\varphi/2)},
\end{equation}
with the horizontal field of view $\varphi = 80^\circ$ ($F = 0.596\,w$). A triangle whose three corners all lie beyond
one side of the view (with a margin of 2\,\%) is left out; one that crosses the near plane $v_z = 6$\,cm is cut there
(a polygon clipped against one plane, then fanned into triangles again).

\paragraph{Filling and depth.} Over a triangle on the screen $1/v_z$ is an affine function of $(s_x, s_y)$, and so are
$u/v_z$ and $v/v_z$ for the texture coordinates; they are interpolated linearly with the edge functions and divided
again per pixel, which is the perspective-correct interpolation. The depth test keeps the larger $1/v_z$, the nearer
surface. The picture is cut into bands of 8 rows that are filled and shaded on several threads.

\paragraph{Light.} Every corner gets a brightness from its normal $\mathbf n$ (turned towards the eye: surfaces are
seen from both sides), the direction to the eye $\mathbf E$ and a fixed light from above, $\mathbf L \propto
(0.35, -0.25, 0.9)$:
\begin{equation}
  b = 0.24 + 0.42\,|\mathbf n\cdot\mathbf L| + 0.3\,\mathbf n\cdot\mathbf E
      + g\bigl(0.7\,(\mathbf n\cdot\mathbf H)^{32} + 0.35\,(\mathbf n\cdot\mathbf E)^{16}\bigr),
  \qquad \mathbf H = \frac{\mathbf L+\mathbf E}{|\mathbf L+\mathbf E|},
\end{equation}
a Lambert term, a term that lights what faces the eye, and for a glossy material ($g>0$) a Blinn highlight and a sheen
towards the eye; the powers are five and four squarings. Per pixel the brightness is the interpolated $b$ times the
material's tone and its procedural texture, plus the emission of a sounding pipe. With the light off (\ui{L}) $b$ is
scaled down to a rest of $0.018$, and every lamp $k$ (a sounding pipe, a flash, a lantern; at most 24, the strongest)
adds
\begin{equation}
  \Delta b = P_k\,\frac{\max(0, \mathbf n\cdot\hat{\boldsymbol\ell}_k)}{1 + d_k^2/R_k^2},
\end{equation}
with the distance $d_k$ and unit direction $\hat{\boldsymbol\ell}_k$ from the pixel to the lamp, its power $P_k$ and
reach $R_k$: an inverse-square fall-off softened within the reach, so a lamp never becomes infinitely bright.

\paragraph{Fog, vignette, grain, the print.} A pixel at depth $d$ is mixed towards the fog's grey $c_f$,
\begin{equation}
  \ell \leftarrow \ell + (c_f-\ell)\,\frac{d/D}{1+d/D},\qquad D = \max\bigl(12,\ 1.5\max(W, L)\bigr)\,\mathrm m,
\end{equation}
half fog at $d = D$, so the far end of any Saal stays readable. Then a vignette $1 - 0.38\,\rho^2/\rho_{\max}^2$
($\rho$ the distance from the middle), a film grain that moves with every frame and is weaker on the highlights,
and a tone curve that gives the greys of an old print,
$g = 0.35\,\ell^2(3-2\ell) + 0.65\,\ell$, $(R,G,B) = (13 + 222g + 6\ell(1-\ell),\ 13 + 219g + 2\ell(1-\ell),\ 12 + 214g -
5\ell(1-\ell))$: a smoothstep for more contrast, the mid-tones a touch warm.

\paragraph{Pipes, fewer sides far away.} A pipe of radius $a$ at depth $d$ is about $\rho = a\,w/(2d)$ pixels wide;
it gets 5, 8, 12 or 18 sides below $\rho = 1.5$, 5, 14 pixels and above. A pipe more than its own length behind the
eye is left out.

\paragraph{Walking.} All the motions relax exponentially: with the frame time $\Delta t$, a quantity $q$ goes towards
its target $q^*$ by
\begin{equation}
  q \leftarrow q + (q^*-q)\,\alpha(\tau),\qquad \alpha(\tau) = 1 - e^{-\Delta t/\tau},
\end{equation}
which is exact for a first-order lag and so does not depend on the frame rate. The mouse sets where to look
($0.22^\circ$ per pixel, $0.15^\circ$ when the pointer is locked), the view follows with $\tau = 35$\,ms; the turning
keys build up a turning speed towards $90^\circ$/s with $\tau = 90$\,ms. The walking speed goes towards $1.6$\,m/s
($4.2$ with Shift) in the direction of the keys with $\tau = 0.12$\,s, dies away with $0.16$\,s when they are let go,
and changes only slowly in the air ($\tau = 0.7$\,s): there is nothing to push against. Against a wall the step is
kept inside the plan (\fn{keepInsidePlan}), and the velocity is then replaced by the step actually made divided by
$\Delta t$, so the part into the wall is gone and the part along it stays, a slide.

A jump starts with $v_z = 4.6$\,m/s against $g = 12.5$\,m/s$^2$ (a little heavier than the earth, so the jump does
not float): the top is $v_z^2/2g = 0.85$\,m above the floor, after $v_z/g = 0.37$\,s, and the feet are back after
$0.74$\,s. Landing, the knees give in by $\min(0.12, 0.016|v_z|)$\,m, which dies away with $\tau = 90$\,ms. In a
loaded model the ground is found by a ray straight down from $0.5$\,m above the feet (the highest step one can climb),
and a step up lifts the feet with $\tau = 50$\,ms instead of at once; a ray upwards keeps the head $0.2$\,m below
what is above it. While walking the eye sways by $0.022\,a\sin\phi$, with the phase growing at $4.4$\,rad per metre
walked and the amount $a\to\min(1, v/1.6)$ (above $0.3$\,m/s, $\tau = 0.15$\,s). The pose goes to the engine with
every frame, and the engine listens from there (\S\ref{sec:walk}): the walked offset moves the listener in the
simulated room, and the yaw turns the head of the headphone decoder and of the AmbiX output.

% ---------------------------------------------------------------------
\subsection{The 360-degree camera (\fn{Pano360.*})}
\label{sec:math-pano}
The button \ui{360\textdegree{} Aufnahme} on the page \ui{Begehen} records a video all around the eye together with
the sound field, for a 360-degree video with spatial sound (YouTube and other players turn both with the viewer's
head). Both are fixed to the room: the picture and the sound are taken as they are before the head's turn, so the
viewer of the video, not the player of the plugin, decides where to look.

\paragraph{The frame's triangles.} While a recording runs, the view's renderer keeps every triangle it is given, in
room coordinates, before any culling (\fn{Raster::setRecording}), and nothing is left out for lying behind the view:
the depth used for leaving out and for the number of sides is then the distance $|\mathbf p - \mathbf e|$ instead of
$v_z$. The triangles of one frame, the materials, the lamps and the glows are written to a file (the triangles of
the room or the model, which come first, only when they differ from the last frame's), and after the recording a
thread of its own draws every frame from it again, from six cameras, the faces of a cube around the eye: front, right, back and left ($\mathbf f = \pm\mathbf x$,
$\pm\mathbf y$), up and down ($\pm\mathbf z$, with $\mathbf r = \mathbf x$), each a square of $N\times N$ pixels with a
field of view of $92^\circ$ (\fn{Raster::setCameraBasis}, which takes the axes directly, as the yaw and pitch of the
view cannot look straight up).

\paragraph{The equirectangular picture.} The video frame of $W\times H$ pixels, $H = W/2$, has the longitude
$\lambda$ along its width and the latitude $\beta$ along its height:
\begin{equation}
  \lambda_i = 2\pi\Bigl(\frac{i+\frac12}{W} - \frac12\Bigr),\qquad
  \beta_j = \pi\Bigl(\frac12 - \frac{j+\frac12}{H}\Bigr),\qquad
  \mathbf d = (\sin\lambda\cos\beta,\ \cos\lambda\cos\beta,\ \sin\beta).
\end{equation}
The middle of the frame is the Saal's front, its right half is the right: a source at azimuth $a$ (positive to the
left) and elevation $e$ appears at $i = W(\frac12 - a/2\pi)$, $j = H(\frac12 - e/\pi)$. The ray $\mathbf d$ meets the
face whose axis is its largest component, $|d_x|$, $|d_y|$ or $|d_z|$; there, with that face's axes,
\begin{equation}
  s_x = \tfrac N2 + F_c\,\frac{\mathbf d\cdot\mathbf r}{\mathbf d\cdot\mathbf f} - \tfrac12,\qquad
  s_y = \tfrac N2 - F_c\,\frac{\mathbf d\cdot\mathbf u}{\mathbf d\cdot\mathbf f} - \tfrac12,\qquad
  F_c = \frac{N/2}{\tan 46^\circ},
\end{equation}
and the pixel is read bilinearly. Since the ray belongs to the face it meets first, $|s_x - N/2|, |s_y - N/2| \le
F_c\tan 45^\circ = 0.966\,N/2$: the extra $2^\circ$ leave a border of $1.7\,\%$ of the face on every side, so the
bilinear read never needs a pixel beyond the edge. The face is made as sharp as the panorama at its equator, where
there are $W/2\pi$ pixels per radian; in the middle of a face there are $F_c$, so
\begin{equation}
  N = \frac{W}{2\pi}\cdot 2\tan 46^\circ = 0.330\,W
\end{equation}
($N = 1266$ for $W = 3840$). Towards the corners of a face the face has more pixels per radian than it needs (up to
$3^{3/2}\approx 5$ times as many per solid angle), so the six faces hold $6N^2 = 0.65\,W^2$ pixels for the $0.5\,W^2$ of
the panorama.

\paragraph{No seams.} Three things of the usual picture depend on where a pixel lies on the screen, and would
show the cube's edges in the panorama. On a face they are made to depend on the ray instead
(\fn{Raster::setPanoFace}):
\begin{itemize}
  \item the fog: the usual picture uses the depth $v_z$; on a face it uses the distance along the ray,
        $d_r = v_z\sqrt{1 + x'^2 + y'^2}$ with $x' = (s_x - N/2)/F_c$, $y' = (N/2 - s_y)/F_c$, the same for both faces
        that meet at an edge;
  \item the dark vault where nothing is (lighter at the horizon): by the ray's elevation $\varepsilon$ instead of the
        row, $c_f\bigl(0.55 + 0.45\max(0, 1 - |\varepsilon|/0.56)\bigr)$;
  \item the vignette: none.
\end{itemize}
The light itself depends only on the surface and the eye (\S\ref{sec:math-walk}), not on the view, so it is the same
on both sides of an edge.

\paragraph{The glows.} The glow of a sounding pipe is drawn over the picture, not into the triangles. Whether a
wall hides it is decided on the face where its middle lies, as in the view (hidden when the depth there is less than
$0.95\,d - 0.4$\,m). It is then laid into the panorama at its longitude and latitude with the angular radius the view
gives it, $\alpha = \tan 40^\circ\,(0.6 + 0.35\,\ell)/d$ for a pipe of length $\ell$ at distance $d$ (limited to
$0.008 \dots 0.34$\,rad), as an ellipse $2\alpha H/\pi$ high and $2\alpha W/(2\pi\cos\beta)$ wide, wider towards
the poles as the equirectangular picture is stretched there, and drawn again one width to the left and the right where
it crosses the seam at the back.

\paragraph{On the graphics card (\fn{PanoGpu.*}).} The drawing thread opens an OpenGL 3.2 context of its own, in a
hidden window, because the view may be closed long before the frames are drawn. For each face the processor still
does what it does for the view's own picture on the card (\S\ref{sec:walkgpu}): it transforms, clips and lights the
triangles; the card fills the face with the view's shader, in which a uniform switches on the three changes above
(the ray $(x', y', 1)$ comes from \texttt{gl\_FragCoord} and the focal length, the vault's elevation from the $z$
components of the face's $\mathbf r$, $\mathbf u$, $\mathbf f$). The six faces are layers of one texture array; a
second pass computes, for every pixel of the panorama, the ray, the face and $(s_x, s_y)$ as above, and reads the
four neighbours with \texttt{texelFetch}, so that the bilinear weights and the rounding are the processor's, not the
card's own filtering. Whether a wall hides a glow is read from the face's depth buffer: the shader writes
$z_w = 1 - 0.06/v_z$, so $v_z = 0.06/(1 - z_w)$. The glows are then laid in by the card as quads with the
sprite's four colour stops, and the picture is read back once, straight into the JPEG's image. Against the
processor's frame the card's differs by 0.03 of 255 on average in the Saal and 0.01 with a model and the light
off (0.08\,\% and 0.002\,\% of the samples by more than 8, along edges where the two rasterise a pixel
differently).

\paragraph{The sound, and the audio as the clock.} The audio thread hands every block's first four Ambisonics
channels to the recorder, taken before the head's turn and before the headphone or loudspeaker decoder: the scene as
the engine renders it, fixed in the room. For ACN / SN3D (AmbiX) the first order of a higher-order scene is exactly
its first four channels, $W, Y, Z, X$; nothing has to be converted, and a player that turns the head by $\psi$
applies the same rotation $R(\psi)$ that the plugin itself applies to its AmbiX output (\fn{AmbiYawRotator}). They
are written as a 24-bit WAV from a lock-free FIFO by a thread of the recorder. Picture and sound are then put
together by the clock of the sound: after $n$ samples at the rate $f_s$ the video frame
\begin{equation}
  k = \Bigl\lfloor \frac{n\,r_v}{f_s} \Bigr\rfloor
\end{equation}
is due ($r_v$ frames per second). The view keeps the scene of the frame it builds whenever a frame $k$ is due that
has no scene yet; that scene then stands for all frames from the first without one up to $k$. A scene is thus never
late by more than one frame of the view (about 17\,ms at its 60 frames per second), and if the view is slower than
$r_v$, frames are repeated, never dropped or shifted: the $k$-th frame of the video always belongs to the moment
$k/r_v$ of the sound, and the video is as long as the sound (at the end the last picture is held to
$\lfloor n_{\text{end}}\,r_v/f_s\rfloor$). Because the drawing comes afterwards, its time does not matter for this:
the recording keeps every frame even at 5760 pixels of width. From the frames and \fn{ambix.wav} ffmpeg makes a MOV
with H.264 and 4 channels of 24-bit PCM at 48\,kHz (sample entry \texttt{lpcm}), the form YouTube asks for.

\paragraph{The marks for YouTube (\fn{SpatialMeta.*}).} A player knows a 360-degree video with spatial sound only by
boxes in the file's \texttt{moov}; the plugin writes the ones Google's Spatial Media Metadata Injector writes
(\texttt{github.com/google/spatial-media}):
\begin{itemize}
  \item Spherical Video V1: a box \texttt{uuid} with the UUID
        \texttt{ffcc8263-f855-\allowbreak 4a93-\allowbreak 8814-\allowbreak 587a02521fdd} and an XML text (\texttt{Spherical}, \texttt{Stitched}, \texttt{ProjectionType} equirectangular), the last child of the
        video's \texttt{trak};
  \item Spherical Video V2, in the video's sample entry (\texttt{avc1}, after its fixed 78 bytes and the codec's own
        boxes, before \texttt{pasp}): \texttt{st3d} (mono) and \texttt{sv3d} with \texttt{svhd}, and \texttt{proj}
        with \texttt{prhd} (yaw, pitch, roll 0) and \texttt{equi} (no crop);
  \item spatial audio: \texttt{SA3D} at the end of the sound's sample entry, 36 bytes -- version 0, periphonic, order 1,
        ACN, SN3D, 4 channels, the channel map $0, 1, 2, 3$ (the file's W, Y, Z, X are already in ACN order). The
        sample entry's fixed part before its boxes depends on its version: 28, 44 or 64 bytes after the header for
        versions 0, 1, 2 (ffmpeg's 4-channel \texttt{lpcm} is version 2, its channel count the 32-bit field at
        byte 40).
\end{itemize}
Only \texttt{moov} changes; the media are copied as they are into a new file. Every box that contains a changed
one grows by as much (\texttt{stsd}, \texttt{stbl}, \texttt{minf}, \texttt{mdia}, \texttt{trak}, \texttt{moov}), and
where \texttt{moov} lies before the media (a \enquote{fast start} file) every chunk offset in \texttt{stco} or
\texttt{co64} that points past it moves by the same $\Delta$. ffprobe reads the marked file as
\enquote{Spherical Mapping: equirectangular} with the sound \enquote{ambisonic 1}, and decodes both the usual and a
fast-start file without an error. Whether YouTube plays it so is the one step not tried yet.
